\chapter{Implementation Code}
\label{appendix:implementations}

This appendix provides complete Python implementations referenced in Chapter~4.

\section{Graph Construction}
\label{appendix:graph_construction}

Heterogeneous graph builder (\texttt{src/features/graph\_builder.py}):

\begin{lstlisting}[language=Python, basicstyle=\ttfamily\fontsize{8}{10}\selectfont]
import polars as pl
import torch
from torch_geometric.data import HeteroData

class HeterogeneousGraphBuilder:
    def __init__(self, entity_types=["device", "ip", "email", "phone", "user"]):
        self.entity_types = entity_types

    def build_graph(self, events_df: pl.DataFrame, cutoff_date: str) -> HeteroData:
        """Build heterogeneous graph with listing nodes and entity-sharing edges."""
        # Filter events before cutoff
        filtered = events_df.filter(pl.col("event_dt") <= cutoff_date)

        # Initialize graph
        graph = HeteroData()

        # Extract node features (44 dims: 36 base + 8 temporal)
        features = self._extract_node_features(filtered)
        graph["listing"].x = torch.tensor(features.to_numpy(), dtype=torch.float)

        # Build edges for each entity type
        for entity_type in self.entity_types:
            edge_index = self._build_sharing_edges(filtered, entity_type)
            graph["listing", f"shares_{entity_type}", "listing"].edge_index = edge_index

        return graph

    def _build_sharing_edges(self, df: pl.DataFrame, entity_type: str) -> torch.Tensor:
        """Create edges between listings sharing the same entity."""
        col_name = f"{entity_type}_hash"

        # Group by entity, get listing IDs
        entity_groups = (
            df.groupby(col_name)
            .agg(pl.col("listing_id").unique())
            .filter(pl.col("listing_id").list.len() > 1)  # Only shared entities
        )

        # Create edges (all pairs within each group)
        edge_list = []
        for group in entity_groups.iter_rows():
            listings = group[1]
            for i, src in enumerate(listings):
                for dst in listings[i+1:]:
                    edge_list.append([src, dst])
                    edge_list.append([dst, src])  # Undirected

        return torch.tensor(edge_list, dtype=torch.long).t()
\end{lstlisting}

Neighbor sampling configuration:

\begin{lstlisting}[language=Python, basicstyle=\ttfamily\fontsize{8}{10}\selectfont]
from torch_geometric.loader import NeighborLoader

# 2-hop neighborhood sampling for training
loader = NeighborLoader(
    graph,
    num_neighbors=[15, 10],  # 15 neighbors at layer 1, 10 at layer 2
    batch_size=512,
    input_nodes=("listing", train_mask),
    shuffle=True,
)
\end{lstlisting}

\section{Hybrid Training Pipeline}
\label{appendix:hybrid_training}

Complete hybrid pipeline (\texttt{src/training/hybrid.py}):

\begin{lstlisting}[language=Python, basicstyle=\ttfamily\fontsize{8}{10}\selectfont]
import numpy as np
import torch
import xgboost as xgb

def train_hybrid_pipeline(train_df, graph, gnn_model, best_params):
    """Train hybrid GNN+XGBoost pipeline."""
    # Stage 1: Generate GNN embeddings
    gnn_model.eval()
    with torch.no_grad():
        embeddings = gnn_model(graph.x_dict, graph.edge_index_dict)
        listing_emb = embeddings["listing"].cpu().numpy()  # [N, 16]

    # Stage 2: Concatenate with tabular features
    tabular_features = train_df[FEATURE_COLUMNS].to_numpy()  # [N, 65]
    combined_features = np.hstack([tabular_features, listing_emb])  # [N, 81]

    # Stage 3: Train XGBoost
    model = xgb.XGBClassifier(**best_params)
    model.fit(
        combined_features,
        train_df["is_fraud"].to_numpy(),
        eval_set=[(X_val, y_val)],
        early_stopping_rounds=10,
        verbose=False
    )

    return model
\end{lstlisting}

\section{Hyperparameter Optimization}
\label{appendix:hpo}

Optuna objective function (\texttt{src/training/hpo.py}):

\begin{lstlisting}[language=Python, basicstyle=\ttfamily\fontsize{8}{10}\selectfont]
import optuna
from optuna.samplers import TPESampler
from sklearn.metrics import average_precision_score

def objective(trial, X_train, y_train, X_val, y_val):
    """Optuna objective for XGBoost hyperparameter optimization."""
    params = {
        "n_estimators": trial.suggest_int("n_estimators", 50, 500, step=50),
        "max_depth": trial.suggest_int("max_depth", 3, 12),
        "learning_rate": trial.suggest_float("learning_rate", 0.01, 0.3, log=True),
        "min_child_weight": trial.suggest_int("min_child_weight", 1, 10),
        "subsample": trial.suggest_float("subsample", 0.6, 1.0),
        "colsample_bytree": trial.suggest_float("colsample_bytree", 0.6, 1.0),
        "scale_pos_weight": trial.suggest_float("scale_pos_weight", 1.0, 20.0),
        "eval_metric": "aucpr",
    }

    model = xgb.XGBClassifier(**params, random_state=42)
    model.fit(
        X_train, y_train,
        eval_set=[(X_val, y_val)],
        early_stopping_rounds=10,
        verbose=False
    )

    y_pred_proba = model.predict_proba(X_val)[:, 1]
    return average_precision_score(y_val, y_pred_proba)

# Run optimization
study = optuna.create_study(direction="maximize", sampler=TPESampler(seed=42))
study.optimize(lambda trial: objective(trial, X_train, y_train, X_val, y_val), n_trials=50)
best_params = study.best_params
\end{lstlisting}

\section{Drift Detection}
\label{appendix:drift_detection}

DDM implementation (\texttt{src/evaluation/drift.py}):

\begin{lstlisting}[language=Python, basicstyle=\ttfamily\fontsize{8}{10}\selectfont]
import numpy as np

class DDM:
    """Drift Detection Method."""
    def __init__(self, warning_threshold=2.0, drift_threshold=3.0):
        self.warning_threshold = warning_threshold
        self.drift_threshold = drift_threshold
        self.errors = []

    def update(self, prediction, true_label):
        """Update with new prediction and return drift status."""
        error = int(prediction != true_label)
        self.errors.append(error)

        # Compute error rate statistics
        p = np.mean(self.errors)
        s = np.sqrt(p * (1 - p) / len(self.errors))

        # Check for drift
        if p + s >= self.drift_threshold:
            self.errors = []  # Reset on drift
            return "DRIFT"
        elif p + s >= self.warning_threshold:
            return "WARNING"
        return "STABLE"
\end{lstlisting}

ADWIN implementation:

\begin{lstlisting}[language=Python, basicstyle=\ttfamily\fontsize{8}{10}\selectfont]
import numpy as np

class ADWIN:
    """Adaptive Windowing drift detector."""
    def __init__(self, delta=0.002):
        self.delta = delta
        self.window = []

    def update(self, prediction, true_label):
        """Update with new prediction and return drift status."""
        error = int(prediction != true_label)
        self.window.append(error)

        # Check all possible splits
        for i in range(1, len(self.window)):
            left = self.window[:i]
            right = self.window[i:]

            if self._detect_change(left, right):
                self.window = self.window[i:]  # Drop old data
                return "DRIFT"

        return "STABLE"

    def _detect_change(self, left, right):
        """Statistical test for distribution change."""
        mu_left = np.mean(left)
        mu_right = np.mean(right)

        n_left = len(left)
        n_right = len(right)
        m = 1 / (1/n_left + 1/n_right)

        var = mu_left * (1 - mu_left) / n_left + mu_right * (1 - mu_right) / n_right
        epsilon_cut = np.sqrt(2 * var * np.log(2 / self.delta) / m)

        return abs(mu_left - mu_right) > epsilon_cut
\end{lstlisting}

\section{Evaluation Metrics}
\label{appendix:metrics}

Metric computation (\texttt{src/utils/metrics.py}):

\begin{lstlisting}[language=Python, basicstyle=\ttfamily\fontsize{8}{10}\selectfont]
from sklearn.metrics import (
    average_precision_score,
    roc_auc_score,
    precision_recall_curve,
    precision_score,
    recall_score,
    f1_score,
)
import numpy as np

def compute_metrics(y_true, y_pred_proba, threshold=0.5):
    """Compute comprehensive evaluation metrics."""
    y_pred = (y_pred_proba >= threshold).astype(int)

    metrics = {
        "auc_pr": average_precision_score(y_true, y_pred_proba),
        "auc_roc": roc_auc_score(y_true, y_pred_proba),
        "precision": precision_score(y_true, y_pred),
        "recall": recall_score(y_true, y_pred),
        "f1": f1_score(y_true, y_pred),
        "precision@50": precision_at_k(y_true, y_pred_proba, k=50),
        "precision@100": precision_at_k(y_true, y_pred_proba, k=100),
        "precision@200": precision_at_k(y_true, y_pred_proba, k=200),
    }

    return metrics

def precision_at_k(y_true, y_scores, k):
    """Compute precision at top-K predictions."""
    sorted_indices = np.argsort(-y_scores)[:k]
    top_k_labels = y_true[sorted_indices]
    return np.sum(top_k_labels) / k
\end{lstlisting}

\section{Feature Engineering}
\label{appendix:feature_engineering}

Cyclical encoding (\texttt{src/features/temporal.py}):

\begin{lstlisting}[language=Python, basicstyle=\ttfamily\fontsize{8}{10}\selectfont]
import numpy as np
import polars as pl

def apply_cyclical_encoding(df: pl.DataFrame) -> pl.DataFrame:
    """Apply cyclical encoding to temporal features."""
    return df.with_columns([
        # Hour encoding (0-23)
        (2 * np.pi * pl.col("event_hour") / 24).sin().alias("hour_sin"),
        (2 * np.pi * pl.col("event_hour") / 24).cos().alias("hour_cos"),
        # Weekday encoding (0-6)
        (2 * np.pi * pl.col("event_day_of_week") / 7).sin().alias("weekday_sin"),
        (2 * np.pi * pl.col("event_day_of_week") / 7).cos().alias("weekday_cos"),
    ])
\end{lstlisting}

Feature schema (\texttt{src/features/schema.py}):

\begin{lstlisting}[language=Python, basicstyle=\ttfamily\fontsize{8}{10}\selectfont]
from dataclasses import dataclass
from typing import List

@dataclass
class FeatureSchema:
    """Feature schema for fraud detection models."""
    seon_boolean: List[str]
    seon_categorical: List[str]
    listing_numeric: List[str]
    listing_categorical: List[str]
    temporal: List[str]

    @property
    def all_features(self) -> List[str]:
        """Return all feature names."""
        return (
            self.seon_boolean
            + self.seon_categorical
            + self.listing_numeric
            + self.listing_categorical
            + self.temporal
        )

    @property
    def categorical_features(self) -> List[str]:
        """Return categorical feature names."""
        return self.seon_categorical + self.listing_categorical

    @property
    def boolean_features(self) -> List[str]:
        """Return boolean feature names."""
        return self.seon_boolean
\end{lstlisting}

\section{ETL Pipeline}
\label{appendix:etl}

As-of join implementation (\texttt{src/data/transform.py}):

\begin{lstlisting}[language=Python, basicstyle=\ttfamily\fontsize{8}{10}\selectfont]
import polars as pl
from typing import Tuple, Dict

def correlate_events_to_seon(
    events_df: pl.DataFrame,
    seon_df: pl.DataFrame,
    seon_unique_col: str = "id",
) -> Tuple[pl.DataFrame, Dict[str, int]]:
    """Correlate each event to closest SEON transaction by timestamp."""

    # Sort both dataframes
    events_sorted = events_df.sort(["INSERTION_ID", "event_dt"])
    seon_sorted = seon_df.sort(["INSERTION_ID", "seon_dt"])

    # Prepare SEON dataframe for join
    seon_for_join = seon_sorted.select([
        "INSERTION_ID", "seon_dt",
        pl.col(seon_unique_col).alias("seon_id"),
        *[col for col in seon_df.columns if col not in ["INSERTION_ID", "seon_dt", seon_unique_col]]
    ])

    # Forward join: closest SEON after the event
    fwd_cols = (
        events_sorted
        .join_asof(
            seon_for_join,
            left_on="event_dt",
            right_on="seon_dt",
            by="INSERTION_ID",
            strategy="forward",
            suffix="_fwd"
        )
    )

    # Backward join: closest SEON before the event
    bwd_cols = (
        events_sorted
        .join_asof(
            seon_for_join,
            left_on="event_dt",
            right_on="seon_dt",
            by="INSERTION_ID",
            strategy="backward",
            suffix="_bwd"
        )
    )

    # Merge and prefer forward match
    merged = fwd_cols.join(
        bwd_cols.select(["INSERTION_ID", "event_dt"] +
                       [col for col in bwd_cols.columns if col.endswith("_bwd")]),
        on=["INSERTION_ID", "event_dt"],
        how="left"
    )

    # Coalesce: prefer forward, fallback to backward
    final_cols = []
    for col in seon_for_join.columns:
        if col not in ["INSERTION_ID", "seon_dt"]:
            final_cols.append(
                pl.coalesce([pl.col(f"{col}_fwd"), pl.col(f"{col}_bwd")]).alias(col)
            )

    result = merged.with_columns(final_cols)

    return result
\end{lstlisting}

\section{MLflow Integration}
\label{appendix:mlflow}

Training with MLflow logging (\texttt{src/training/trainer.py}):

\begin{lstlisting}[language=Python, basicstyle=\ttfamily\fontsize{8}{10}\selectfont]
import mlflow
from omegaconf import OmegaConf

def train_with_mlflow(cfg, pipeline):
    """Train model with MLflow experiment tracking."""

    # Configure MLflow
    mlflow.set_tracking_uri(cfg.experiment.tracking_uri)
    mlflow.set_experiment(cfg.experiment.name)

    # Start run
    run_name = f"{cfg.model.variant}_{cfg.data.train_start_date}"
    with mlflow.start_run(run_name=run_name):

        # Log hyperparameters
        params_dict = OmegaConf.to_container(cfg, resolve=True)
        mlflow.log_params(flatten_dict(params_dict))

        # Train model
        result = pipeline.run()

        # Log metrics
        mlflow.log_metrics({
            "train_auc_pr": result.train_metrics.auc_pr,
            "train_auc_roc": result.train_metrics.auc_roc,
            "val_auc_pr": result.val_metrics.auc_pr,
            "test_auc_pr": result.test_metrics.auc_pr,
            "test_precision@50": result.test_metrics.precision_at_50,
            "test_precision@100": result.test_metrics.precision_at_100,
            "training_time_seconds": result.training_time,
        })

        # Log model artifact
        mlflow.xgboost.log_model(result.model, "model")

        # Log SHAP plots if enabled
        if cfg.training.run_shap:
            mlflow.log_artifact(result.shap_summary_plot, "shap_summary.png")

        return result

def flatten_dict(d, parent_key='', sep='_'):
    """Flatten nested dictionary for MLflow logging."""
    items = []
    for k, v in d.items():
        new_key = f"{parent_key}{sep}{k}" if parent_key else k
        if isinstance(v, dict):
            items.extend(flatten_dict(v, new_key, sep=sep).items())
        else:
            items.append((new_key, v))
    return dict(items)
\end{lstlisting}

\section{API Service}
\label{appendix:api}

FastAPI endpoint (\texttt{src/api/main.py}):

\begin{lstlisting}[language=Python, basicstyle=\ttfamily\fontsize{8}{10}\selectfont]
from fastapi import FastAPI, HTTPException
from pydantic import BaseModel
import xgboost as xgb
import shap

app = FastAPI(title="Fraud Detection API")

# Load model at startup
model = xgb.Booster()
model.load_model("models/production_model.json")
explainer = shap.TreeExplainer(model)

class EventPayload(BaseModel):
    insertion_id: str
    device_hash: str
    ip_hash: str
    email_hash: str
    # ... other features

class PredictResponse(BaseModel):
    probability: float
    risk_level: str
    top_factors: list

@app.post("/predict")
async def predict(event: EventPayload) -> PredictResponse:
    """Generate fraud prediction with explanation."""
    try:
        # Extract features
        features = feature_engineer.extract(event, event_store)

        # Enrich with graph-like counts
        features["device_link_count"] = event_store.count_related("device", event.device_hash)
        features["ip_link_count"] = event_store.count_related("ip", event.ip_hash)
        features["email_link_count"] = event_store.count_related("email", event.email_hash)

        # Generate prediction
        dmatrix = xgb.DMatrix([features])
        prob = model.predict(dmatrix)[0]

        # Compute SHAP values
        shap_values = explainer.shap_values(dmatrix)
        top_factors = get_top_factors(shap_values[0], feature_names, k=5)

        # Determine risk level
        if prob > 0.7:
            risk_level = "HIGH"
        elif prob > 0.3:
            risk_level = "MEDIUM"
        else:
            risk_level = "LOW"

        return PredictResponse(
            probability=float(prob),
            risk_level=risk_level,
            top_factors=top_factors
        )

    except Exception as e:
        raise HTTPException(status_code=500, detail=str(e))

def get_top_factors(shap_values, feature_names, k=5):
    """Get top-K contributing factors."""
    importance = [(name, abs(val)) for name, val in zip(feature_names, shap_values)]
    importance.sort(key=lambda x: x[1], reverse=True)
    return [{"feature": name, "importance": float(val)} for name, val in importance[:k]]
\end{lstlisting}

\section{GNN Implementation}
\label{appendix:gnn}

GraphSAGE encoder (\texttt{src/models/graphsage.py}):

\begin{lstlisting}[language=Python, basicstyle=\ttfamily\fontsize{8}{10}\selectfont]
import torch
import torch.nn as nn
import torch.nn.functional as F
from torch_geometric.nn import SAGEConv, HeteroConv

class GraphSAGEEncoder(torch.nn.Module):
    """GraphSAGE encoder for heterogeneous graphs."""

    def __init__(self, in_channels, hidden_channels, out_channels,
                 num_layers, dropout, edge_types):
        super().__init__()
        self.convs = nn.ModuleList()

        for i in range(num_layers):
            in_ch = in_channels if i == 0 else hidden_channels
            out_ch = out_channels if i == num_layers - 1 else hidden_channels

            # Create SAGEConv for each edge type
            conv_dict = {
                ("listing", edge_type, "listing"): SAGEConv(in_ch, out_ch, aggr="mean")
                for edge_type in edge_types
            }
            self.convs.append(HeteroConv(conv_dict, aggr="sum"))

        self.dropout = dropout

    def forward(self, x_dict, edge_index_dict):
        """Forward pass through GNN layers."""
        x = x_dict
        for i, conv in enumerate(self.convs):
            x = conv(x, edge_index_dict)
            if i < len(self.convs) - 1:  # No activation/dropout on last layer
                x = {k: F.relu(v) for k, v in x.items()}
                x = {k: F.dropout(v, p=self.dropout, training=self.training)
                     for k, v in x.items()}

        return x["listing"]

def train_supervised(model, loader, optimizer, device, epochs=30):
    """Supervised GNN training loop."""
    model.train()
    for epoch in range(epochs):
        total_loss = 0
        for batch in loader:
            batch = batch.to(device)
            optimizer.zero_grad()

            # Forward pass
            out = model(batch.x_dict, batch.edge_index_dict)

            # Compute loss on training nodes
            loss = F.binary_cross_entropy_with_logits(
                out[batch["listing"].train_mask],
                batch["listing"].y[batch["listing"].train_mask],
                pos_weight=torch.tensor([11.0], device=device)  # Class imbalance
            )

            # Backward pass
            loss.backward()
            optimizer.step()

            total_loss += loss.item()

        avg_loss = total_loss / len(loader)
        print(f"Epoch {epoch+1}/{epochs}, Loss: {avg_loss:.4f}")
\end{lstlisting}
